\documentclass[11pt]{article}
\usepackage[a4paper,margin=28mm]{geometry}
\usepackage[T1]{fontenc}
\usepackage{lmodern}
\usepackage{amsmath,amssymb,amsthm}
\usepackage[hidelinks]{hyperref}
\newtheorem{theorem}{Theorem}
\title{A lower bound for the small-time cost of boundary control of the one-dimensional Schr\"odinger equation}
\author{Kaiwen Jin\\\normalsize\href{mailto:yui.ui.siki@gmail.com}{\texttt{yui.ui.siki@gmail.com}}}
\date{}
\begin{document}
\maketitle
\begin{abstract}
We consider the one-dimensional Schr\"odinger equation on an interval of
length $L$, driven by an $L^2$ Dirichlet control at one endpoint, with initial
data measured in $H^{-1}$. We prove that its null-control norm cost satisfies
$C_S(T,L)\ge c_L T^{7/2}\exp(L^2/(\pi T))$ for all sufficiently small $T$.
Thus $\liminf_{T\downarrow0}T\log C_S(T,L)\ge L^2/\pi$, excluding the proposed
coefficient $L^2/4$. The proof uses one sine initial state: its exact control
moments prescribe the zeros of an entire function, and Jensen's formula turns
those zeros into the lower bound.
\end{abstract}

\section{Introduction and main result}\label{sec:intro}
Boundary controllability asks how strongly a system must be driven at its
boundary to reach a prescribed final state. For a Schr\"odinger equation on a
segment, the small-time question concerns the growth of the required control
norm as the available time tends to zero. It is closely related to the
quantitative moment problem and entire-function methods developed for this
setting; see Miller~\cite[Section 4]{Miller2003}.

Fix $L>0$. We consider complex-valued transposition solutions of
\begin{equation}\label{eq:pde}
 \begin{cases}
 i y_t+y_{xx}=0,&(t,x)\in(0,T)\times(0,L),\\
 y(t,0)=h(t),\quad y(t,L)=0,&0<t<T,\\
 y(0,\cdot)=y_0\in H^{-1}(0,L),&
 \end{cases}
\end{equation}
where $h\in L^2(0,T;\mathbb C)$. We use $H^{-1}$ as the continuous anti-dual
of $H_0^1$, with a fixed standard Sobolev norm independent of $T$.
The proof applies both to the derivative norm on $H_0^1$ and to its usual
full Sobolev norm. Define the norm control cost by
\begin{equation}\label{eq:cost}
 C_S(T,L)=\sup_{\|y_0\|_{H^{-1}}=1}\,
 \inf\bigl\{\|h\|_{L^2(0,T)}:\ y(0)=y_0,\ y(T)=0
 \text{ in }\eqref{eq:pde}\bigr\}.
\end{equation}
The infimum of an empty set and $\log(+\infty)$ are understood in the extended
sense. In particular, $C_S$ is a norm cost, rather than its square.

The lower coefficient $L^2/4$ appears in
Lissy~\cite[Theorem 1.2, equation (8)]{Lissy2015}.
For precisely the norm-cost setting~\eqref{eq:pde}--\eqref{eq:cost},
Lissy~\cite[Sections 3.1--3.2]{Lissy2025} conjectures
\begin{equation}\label{eq:question}
 \lim_{T\downarrow0}T\log C_S(T,L)=\frac{L^2}{4}.
\end{equation}
The same work gives the upper bound
$\limsup_{T\downarrow0}T\log C_S(T,L)\le 27^{1/4}L^2/4$
\cite[Theorem 3.1]{Lissy2025}. Here we prove a strict improvement of the lower
coefficient, which rules out~\eqref{eq:question}.

\begin{theorem}\label{thm:main}
For every fixed $L>0$, there is a constant $c_L>0$, independent of $T$, such
that
\begin{equation}\label{eq:main-bound}
 C_S(T,L)\ge c_L T^{7/2}\exp\!\left(\frac{L^2}{\pi T}\right),
 \qquad 0<T\le\frac{L^2}{2\pi}.
\end{equation}
Consequently,
\begin{equation}\label{eq:liminf}
 \liminf_{T\downarrow0}T\log C_S(T,L)
 \ge\frac{L^2}{\pi}>\frac{L^2}{4}.
\end{equation}
\end{theorem}

The constant $c_L$ may depend on the fixed interval and Sobolev norm. The
statement covers the full unit sphere in~\eqref{eq:cost}; a single unit sine
state supplies its lower bound. It does not determine the sharp coefficient
or establish that the limit exists.

The argument uses only necessary control moments. After shifting the first
spectral frequency to zero, a centered transform of each control takes the
value one at the origin and vanishes at all remaining quadratic spectral
gaps. Its growth on a circle is bounded by the control norm and the time
support. Jensen's formula compares these two facts. An exact factorial product
then gives the exponential lower bound. Sections~\ref{sec:moments}
and~\ref{sec:jensen} derive the moments and the zero estimate;
Section~\ref{sec:bound} chooses the radius and completes the proof.

\section{Necessary moments for one unit initial state}\label{sec:moments}
We use the pairing linear in the state and conjugate-linear in the test
function. For a transposition solution and a smooth homogeneous Dirichlet
adjoint solution $z$, the endpoint identity is
\begin{equation}\label{eq:tp}
 \langle y(T),z(T)\rangle-\langle y_0,z(0)\rangle
 =i\int_0^T h(t)\overline{z_x(t,0)}\,dt.
\end{equation}
Here the adjoint solves $iz_t+z_{xx}=0$ backwards from its terminal datum.
For smooth states, differentiating $\int_0^L y\overline z\,dx$ gives
$i[y_x\overline z-y\overline{z_x}]_0^L
=i h\overline{z_x(t,0)}$.
Identity~\eqref{eq:tp} is its transposed form; it requires no spatial trace
of a rough state.

Let
\begin{equation}\label{eq:spectrum}
 \phi_n(x)=\sqrt{\frac2L}\sin\frac{n\pi x}{L},\qquad
 \lambda_n=\frac{n^2\pi^2}{L^2},\qquad
 b_n=\phi_n'(0)=\sqrt{\frac2L}\frac{n\pi}{L}.
\end{equation}
These functions are orthonormal in $L^2(0,L)$ and satisfy
$-\phi_n''=\lambda_n\phi_n$ with zero endpoint values. Set
\begin{equation}\label{eq:witness}
 a_L=\|\phi_1\|_{H^{-1}}^{-1},\qquad
 y_0=a_L\phi_1,\qquad d_L=\frac{a_L}{b_1}>0.
\end{equation}
The norm in this definition is finite and nonzero: a smooth sine defines a
continuous functional on $H_0^1$, and its pairing with itself is one.
Thus $y_0$ has exactly unit $H^{-1}$ norm. These quantities depend on $L$
and the fixed Sobolev norm, but not on $T$.

For every control taking this state to zero, test~\eqref{eq:tp} with
$z(t,x)=e^{-i\lambda_n(t-T)}\phi_n(x)$. Since
$z(0)=e^{i\lambda_nT}\phi_n$, we obtain
\[
 -a_L\delta_{n1}e^{-i\lambda_nT}
 =i b_n e^{-i\lambda_nT}\int_0^T h(t)e^{i\lambda_n t}\,dt.
\]
Consequently every such control satisfies
\begin{equation}\label{eq:moments}
 \int_0^T h(t)e^{i\lambda_n t}\,dt=i d_L\delta_{n1},\qquad n\ge1.
\end{equation}
Only smooth single-mode tests have been used. In particular, this necessary
condition does not involve an infinite expansion of the controlled state.

\section{The entire transform and its zeros}\label{sec:jensen}
Normalize the nonzero moment and center the time interval by defining
\begin{equation}\label{eq:transform}
 v(t)=\frac{h(t)e^{i\lambda_1t}}{i d_L},\qquad
 F(z)=\int_0^T v(t)e^{iz(t-T/2)}\,dt,\qquad z\in\mathbb C.
\end{equation}
Cauchy--Schwarz gives $\|v\|_1\le\sqrt T\,\|v\|_2
=\sqrt T\,\|h\|_2/d_L$. The transform is entire. Indeed, on $|z|\le A$,
its exponential series can be integrated term by term, with remainder at
most
\[
 \|v\|_1\sum_{k>K}\frac{(AT/2)^k}{k!}\longrightarrow0.
\]
The resulting power series has infinite radius of convergence. Also,
\begin{equation}\label{eq:growth}
 |F(z)|\le K_T e^{(T/2)|\operatorname{Im}z|},\qquad
 K_T=\frac{\sqrt T\,\|h\|_2}{d_L}.
\end{equation}
The factor $T/2$ is the reason for centering the interval.

The moments~\eqref{eq:moments} imply
\begin{equation}\label{eq:zeros}
 F(0)=1,\qquad F(\mu_n)=0\quad(n\ge2),\qquad
 \mu_n=\lambda_n-\lambda_1=\frac{\pi^2}{L^2}(n^2-1).
\end{equation}
Thus $F$ is nonzero and has distinct positive prescribed zeros. It may have
other zeros, which only strengthen the estimate below.

Let $r>0$ be a radius with no zero on its circle. Jensen's formula at zero,
with zeros $z_j$ counted according to multiplicity, reads
\[
 \sum_{|z_j|<r}\log\frac r{|z_j|}
 =\frac1{2\pi}\int_0^{2\pi}\log|F(re^{i\theta})|\,d\theta,
\]
because $F(0)=1$. The sum is finite, since zeros of a nonzero entire function
are isolated. Each summand is nonnegative. Taking only the prescribed zeros
and applying~\eqref{eq:growth}, we find
\[
 \sum_{\mu_n<r}\log\frac r{\mu_n}
 \le\log K_T+\frac{Tr}{\pi},
\]
where $\int_0^{2\pi}|\sin\theta|\,d\theta=4$ fixes the constant.
For any $R>\mu_N$, with an integer $N\ge2$, it follows that
\begin{equation}\label{eq:jensen}
 \log K_T\ge\sum_{n=2}^{N}\log\frac R{\mu_n}-\frac{TR}{\pi}.
\end{equation}
To include radii whose circles contain additional zeros, choose zero-free
radii increasing to $R$, all larger than $\mu_N$. Such radii exist because
the zero set is countable. Pass to the limit only in the selected finite
sum. This proves~\eqref{eq:jensen} at every required $R$.

\section{The lower bound}\label{sec:bound}
Choose $R=\pi^2N^2/L^2$, so $\mu_N<R$. The prescribed-zero contribution is
\begin{align}
 S_N:=\sum_{n=2}^{N}\log\frac R{\mu_n}
 &=2(N-1)\log N-\log\frac{(N-1)!(N+1)!}{2}.\label{eq:product}
\end{align}
Here we used $\prod_{n=2}^{N}(n^2-1)=(N-1)!(N+1)!/2$.
Since $\log$ is increasing,
\[
 \log N!\le\int_1^N\log x\,dx+\log N
 =N\log N-N+1+\log N.
\]
Together with $(N-1)!(N+1)!=(N!)^2(1+1/N)$, this gives, for every $N\ge2$,
\begin{equation}\label{eq:sum-bound}
 S_N\ge2N-4\log N-2+\log2-\log(1+1/N)
 \ge2N-4\log N-2.
\end{equation}
Substituting into~\eqref{eq:jensen} and recalling~\eqref{eq:growth}, we obtain
\begin{equation}\label{eq:integer-bound}
 \|h\|_2\ge\frac{d_L}{\sqrt T}\,N^{-4}e^{-2}
 \exp\!\left(2N-\frac{\pi T}{L^2}N^2\right),\qquad N\ge2.
\end{equation}
This holds for every admissible control of the fixed unit state and every
integer $N\ge2$.

\begin{proof}[Proof of Theorem~\ref{thm:main}]
For $0<T\le L^2/(2\pi)$, set
$X=L^2/(\pi T)\ge2$, $N=\lfloor X\rfloor$ and $a=1/X\le1/2$.
Completing the square gives
\[
 2N-aN^2=X-a(N-X)^2\ge X-\tfrac12,
 \qquad N^{-4}\ge X^{-4}=\left(\frac{\pi T}{L^2}\right)^4.
\]
Therefore~\eqref{eq:integer-bound} implies
\begin{equation}\label{eq:control-bound}
 \|h\|_2\ge c_L T^{7/2}e^{L^2/(\pi T)},\qquad
 c_L=d_L\pi^4L^{-8}e^{-5/2}>0.
\end{equation}
The constant is independent of $T$ and of the control. Taking the infimum
over all controls for the unit state~\eqref{eq:witness}, and then the
supremum over the whole unit sphere, proves~\eqref{eq:main-bound}.
Taking the logarithm, multiplying by $T$, and using $T\log T\to0$ proves
\eqref{eq:liminf}. The strict inequality follows from $\pi<4$.
\end{proof}

\paragraph{AI-use statement.}
OpenAI Codex was used for literature searches, proof development, calculation
checks, and drafting this manuscript. The author is responsible for its
mathematical claims.

\begin{thebibliography}{9}
\bibitem{Miller2003}
Luc Miller,
\emph{How violent are fast controls for Schr\"odinger and plate vibrations?},
arXiv:math/0309040v3, 2003.
\url{https://arxiv.org/abs/math/0309040v3}.

\bibitem{Lissy2015}
Pierre Lissy,
\emph{Explicit lower bounds for the cost of fast controls for some 1-D
parabolic or dispersive equations, and a new lower bound concerning the
uniform controllability of the 1-D transport-diffusion equation},
author manuscript, 8 June 2015.
\url{https://cermics.gitlab.enpc.fr/lissy/Publi/chal.pdf}.

\bibitem{Lissy2025}
Pierre Lissy,
\emph{Effective multipliers for weights whose log are H\"older continuous.
Application to the cost of fast boundary controls for the 1D Schr\"odinger
equation}, arXiv:2502.04859v1, 2025.
\url{https://arxiv.org/abs/2502.04859v1}.
\end{thebibliography}
\end{document}
